\label{chap:pantallas-dualidad}
\index{pantalla algebraica}
\index{dualidad residuo--cociente}

Los enteros \(10\), \(11\), \(24\) y \(26\) aparecen en la construcción por
procedimientos algebraicos diferentes: longitud de un código, dimensión de un
módulo, rango de un retículo y dimensión de un cociente ortogonal. El
propósito de esta sección es reconstruir esos procedimientos y los mapas que
los relacionan. Una igualdad entre enteros no bastará para identificar los
objetos que éstos numeran.

Llamaremos \emph{pantalla algebraica} a un subespacio o cociente acompañado
del mapa que lo selecciona. Esta denominación no introduce una nueva clase de
espacios: abrevia el par formado por el objeto lineal y su proyector o
aplicación cociente. Con esta convención, las reducciones

\[
12\longrightarrow11\longrightarrow10,
\qquad
26\longrightarrow24
\]

son enunciados sobre aplicaciones definidas, no descomposiciones numerológicas.
La primera cadena procede de una representación de \(A_5\) marcada por el
vector dodecafásico \(K\); la segunda, de una reducción isotrópica de un
retículo lorentziano. Ambas se apoyan en la rama excepcional construida en
las secciones precedentes, pero no son el mismo mecanismo.

\section{El cociente perfecto de \(11\) coordenadas}
\label{sec:golay-once}

Sea

\[
C_W=[12,6,6]_3
\]

el código ternario extendido de Golay. La perforación de una coordenada
produce el código ternario de Golay

\[
G_{11}=[11,6,5]_3\subset\mathbb F_3^{11}
\]

\cite{MacWilliamsSloane1977}. La longitud once pertenece aquí al espacio
ambiente del código; no se introduce mediante una interpretación física.

\begin{definition}[Bola ternaria]
Para \(x\in\mathbb F_3^n\) y \(t\geq0\), la bola de Hamming de radio \(t\)
centrada en \(x\) es
\[
B_t(x)=\{y\in\mathbb F_3^n:d_H(x,y)\leq t\}.
\]
Su cardinal, independiente de \(x\), es
\[
V_3(n,t)=\sum_{j=0}^{t}\binom nj2^j.
\]
\end{definition}

\begin{theorem}[Perfección de \(G_{11}\)]
\label{thm:g11-perfecto}
Las bolas de Hamming de radio dos centradas en las palabras de \(G_{11}\)
constituyen una partición de \(\mathbb F_3^{11}\). En particular,
\[
V_3(11,2)=243,
\qquad
|G_{11}|\,V_3(11,2)=3^{11}.
\]
\end{theorem}

\begin{proof}
La distancia mínima de \(G_{11}\) es cinco. Si dos bolas de radio dos,
centradas en palabras distintas \(c,c'\in G_{11}\), tuvieran un elemento
común \(y\), la desigualdad triangular daría
\[
d_H(c,c')\leq d_H(c,y)+d_H(y,c')\leq4,
\]
en contradicción con \(d_H(c,c')\geq5\). Las bolas son, por tanto,
disjuntas. Su cardinal es
\[
V_3(11,2)
=\binom{11}{0}+2\binom{11}{1}+2^2\binom{11}{2}
=1+22+220=243=3^5.
\]
Como \(G_{11}\) tiene dimensión seis, contiene \(3^6\) palabras. En
consecuencia,
\[
3^6V_3(11,2)=3^6\cdot3^5=3^{11}
=|\mathbb F_3^{11}|.
\]
La unión disjunta de las bolas tiene el cardinal del espacio ambiente y, por
tanto, coincide con él.
\end{proof}

La decodificación perfecta asigna a cada palabra recibida una única palabra
de código a distancia a lo sumo dos. Equivalentemente, el espacio de
síndromes tiene \(3^{11-6}=243\) elementos. Esta aparición de \(243\) debe
distinguirse de la capa cúbica

\[
3\cdot9^2=243
\]

de la completación \(1000=729+243+27+1\). Una cantidad enumera síndromes de
un código; la otra enumera posiciones de una capa de frontera. La igualdad
de cardinales autoriza una comparación, pero no proporciona una biyección
canónica entre ambas realizaciones.

También conviene separar la partición métrica

\[
243=1+22+220
\]

de la identidad tipada

\[
243=1+22+90+120+10.
\]

En la segunda expresión, \(90=6\binom62\) y
\(120=8\binom62\) cuentan familias de incidencias en la extensión
hexada--octada, mientras que \(10\) será la dimensión de una pantalla lineal.
No son subconjuntos de pesos \(0,1,2\) de una misma bola de Hamming.
La construcción de diseños que sustenta esa extensión se inscribe en la
teoría clásica de los sistemas de Witt \cite{AssmusKey1992}.

\section{\(2\) descendientes del código ternario extendido}
\label{sec:dos-descendientes-cw}

La longitud once no es el único rango que desciende de \(C_W\). El mismo
código interviene en el pegado del retículo de raíces \(A_2^{12}\), cuyo
rango es veinticuatro. Después de fijar los datos de cámara y el vector de
vecindad construidos en la rama excepcional, el vecino de índice tres sin
raíces es el retículo de Leech \(\Lambda_{24}\). La estructura de
dependencias es

\begin{equation}
\label{eq:dos-descendientes-cw}
\begin{gathered}
C_W\xrightarrow{\text{perforación}}G_{11}\subset\mathbb F_3^{11},\\[2mm]
C_W\xrightarrow{\text{pegado}}N(A_2^{12})
\xrightarrow{\text{vecino de índice }3}\Lambda_{24}.
\end{gathered}
\end{equation}

La primera flecha es un morfismo de códigos; la segunda línea pertenece a la
teoría de retículos. La fuente ternaria común explica la coexistencia de los
rangos once y veinticuatro sin identificarlos. En particular, \(11\) no es
una subdimensión de \(24\), ni \(24\) se obtiene completando el código
perforado.

\begin{proposition}[Bifurcación de rangos]
La construcción \eqref{eq:dos-descendientes-cw} produce dos invariantes de
tipo distinto:
\[
\operatorname{length}(G_{11})=11,
\qquad
\operatorname{rank}(\Lambda_{24})=24.
\]
Su procedencia común es el código \(C_W\); ninguna de las dos igualdades
implica un isomorfismo entre el espacio de coordenadas del código y el
espacio real generado por el retículo.
\end{proposition}

\begin{proof}
La perforación elimina una de las doce posiciones y no altera la dimensión
seis del código, de donde la longitud once. Cada copia de \(A_2\) tiene rango
dos; por ello \(A_2^{12}\), sus pegados de índice finito y sus vecinos de
rango completo tienen rango veinticuatro. Los objetos viven, respectivamente,
en las categorías de códigos sobre \(\mathbb F_3\) y retículos enteros. Un
isomorfismo entre ellos ni siquiera estaría bien tipado sin un funtor
adicional.
\end{proof}

\section{La representación icosaédrica de dimensión \(12\)}
\label{sec:representacion-a5-doce}

La reducción \(12\to11\to10\) procede de una segunda realización del número
doce. Sea \(V_{12}=\mathbb R^{12}\) el módulo de permutación de \(A_5\) sobre
los doce vértices de un icosaedro. Denotemos por \(\rho\) esta representación
y por \(\chi_{12}\) su carácter. Para fijar sin ambigüedad la correspondencia
con las doce posiciones dodecafásicas, realizamos los vértices como las clases laterales
izquierdos \(A_5/H\), donde \(A_5\) actúa sobre \(\{0,1,2,3,4\}\) y
\[
 H=\langle(0\ 1\ 2\ 3\ 4)\rangle\cong C_5.
\]
La posición \(j\) se identifica con la clase lateral \(r_jH\), para los
representantes siguientes, escritos en notación de imágenes:
\[
\begin{array}{c|c@{\qquad}c|c@{\qquad}c|c}
j&r_j&j&r_j&j&r_j\\ \hline
1&(3,2,4,1,0)&5&(3,2,0,4,1)&9&(3,1,2,4,0)\\
2&(3,1,0,2,4)&6&(4,0,1,2,3)&10&(2,1,4,3,0)\\
3&(0,1,3,4,2)&7&(4,3,2,1,0)&11&(3,4,1,2,0)\\
4&(1,4,2,3,0)&8&(0,2,3,1,4)&12&(4,2,1,3,0)
\end{array}
\]
Esta tabla es el dato de orientación que permite aplicar un proyector de
\(A_5\) al vector \(K\); sin ella, una norma coordenada dependería de una
permutación no especificada. En el orden de clases conjugadas

\[
1A,\quad2A,\quad3A,\quad5A,\quad5B,
\]

los tamaños de las clases y el número de puntos fijos son

\[
\begin{array}{c|rrrrr}
 &1A&2A&3A&5A&5B\\ \hline
|C|&1&15&20&12&12\\
\chi_{12}&12&0&0&2&2.
\end{array}
\]

En efecto, las rotaciones de órdenes dos y tres no fijan vértices, mientras
que cada rotación de orden cinco fija los dos vértices de su eje. Escribamos

\[
\phi=\frac{1+\sqrt5}{2},
\qquad
\phi'=\frac{1-\sqrt5}{2},
\qquad
\phi+\phi'=1.
\]

La parte relevante de la tabla de caracteres de \(A_5\) es

\[
\begin{array}{c|rrrrr}
 &1A&2A&3A&5A&5B\\ \hline
\chi_1&1&1&1&1&1\\
\chi_3&3&-1&0&\phi&\phi'\\
\chi_{3'}&3&-1&0&\phi'&\phi\\
\chi_4&4&0&1&-1&-1\\
\chi_5&5&1&-1&0&0.
\end{array}
\]

\begin{theorem}[Descomposición del módulo icosaédrico]
\label{thm:descomposicion-a5}
La representación de permutación sobre los doce vértices se descompone como
\[
V_{12}\cong V_1\oplus V_3\oplus V_{3'}\oplus V_5.
\]
El subespacio ortogonal al vector uniforme es
\[
V_{11}:=\mathbf1^\perp
\cong V_3\oplus V_{3'}\oplus V_5
\]
y tiene dimensión once.
\end{theorem}

\begin{proof}
La multiplicidad de un carácter irreducible \(\chi\) en \(\chi_{12}\) es
\[
\langle\chi_{12},\chi\rangle
=\frac1{60}\sum_C|C|\chi_{12}(C)\chi(C).
\]
Sólo contribuyen \(1A,5A,5B\). Para los caracteres de dimensiones uno, tres,
tres prima, cuatro y cinco se obtienen, respectivamente,
\[
\frac{12+24+24}{60}=1,
\]
\[
\frac{36+24(\phi+\phi')}{60}=1,
\qquad
\frac{36+24(\phi'+\phi)}{60}=1,
\]
\[
\frac{48-24-24}{60}=0,
\qquad
\frac{60}{60}=1.
\]
Por tanto, las multiplicidades son \((1,1,1,0,1)\). El único vector fijo por
todas las permutaciones es el vector uniforme
\(\mathbf1=(1,\ldots,1)\). Su complemento ortogonal contiene los tres
sumandos no triviales y tiene dimensión \(3+3+5=11\).
\end{proof}

El proyector ortogonal sobre \(V_{11}\) es explícito:

\[
P_{11}=I_{12}-\frac1{12}J_{12},
\]

donde \(J_{12}\) es la matriz cuyas entradas son uno. Así, la primera
reducción es la aplicación

\[
\mathbb R^{12}\longrightarrow V_{11},
\qquad
x\longmapsto x-\frac1{12}\Bigl(\sum_{j=1}^{12}x_j\Bigr)\mathbf1.
\]

Esta fórmula especifica qué información se elimina: exclusivamente la
componente uniforme. La igualdad \(12=1+11\) es una consecuencia de la
descomposición del módulo, no una convención dimensional.

\section[Polarización del sector tridimensional]{Polarización del sector tridimensional determinada por K}
\label{sec:reduccion-once-diez}

Consideremos el vector dodecafásico ya obtenido,

\[
K=(234,543,140,729,659,824,621,58,914,794,146,601),
\qquad
\sum_{j=1}^{12}K_j=6263.
\]

Su componente de suma cero es

\[
\widetilde K=P_{11}K
=K-\frac{6263}{12}\mathbf1.
\]

Para seleccionar una de las dos representaciones tridimensionales usamos la
bijección posición--clase lateral fijada en la tabla anterior. El proyector central
correspondiente se escribe

\[
P_{\mathrm{ico},3}=\frac{3}{60}\sum_{g\in A_5}\chi_3(g^{-1})\rho(g).
\]

La suma contiene matrices de permutación y coeficientes en
\(\mathbb Q(\sqrt5)\); por consiguiente, \(P_{\mathrm{ico},3}\) es un operador exacto sobre
ese cuerpo, satisface \(P_{\mathrm{ico},3}^2=P_{\mathrm{ico},3}=P_{\mathrm{ico},3}^{\mathsf T}\) y tiene rango tres.

\begin{definition}[Dirección polarizante]
La proyección icosaédrica del vector y la recta que genera son
\[
u_K=P_{\mathrm{ico},3}\widetilde K,
\qquad
\ell_K=\mathbb R u_K.
\]
\end{definition}

\begin{lemma}[No degeneración de la polarización]
\label{lem:norma-uk}
Para el vector \(K\) y la correspondencia posición--clase lateral declarada,
\[
\|u_K\|^2
=\frac{6638585+2275584\sqrt5}{20}>0.
\]
En particular, \(\ell_K\) es una recta bien definida.
\end{lemma}

\begin{proof}
Al sustituir la expresión de \(P_{\mathrm{ico},3}\) en
\(\|u_K\|^2=\widetilde K^{\mathsf T}P_{\mathrm{ico},3}\widetilde K\), la suma se reduce a
\[
\frac3{60}\sum_{g\in A_5}
\chi_3(g^{-1})\,
\widetilde K^{\mathsf T}\rho(g)\widetilde K.
\]
Cada producto escalar es racional, pues \(\rho(g)\) es una matriz de
permutación y \(\widetilde K\in\mathbb Q^{12}\). La separación de las dos
clases de orden cinco deja como única parte irracional un múltiplo de
\(\sqrt5\). La suma exacta de los sesenta términos es
\[
\frac{6638585}{20}+\frac{2275584}{20}\sqrt5.
\]
Ambos coeficientes son positivos y \(\sqrt5>0\); por tanto, la norma no se
anula.
\end{proof}

La enumeración exacta del grupo y de las doce clases laterales reconstruye el
proyector sobre \(\mathbb Q(\sqrt5)\), comprueba su idempotencia y produce la
norma anterior.

\begin{theorem}[Reducción polarizada \(11\to10\)]
\label{thm:once-diez}
Sea
\[
V_{10}:=(V_3\cap\ell_K^\perp)\oplus V_{3'}\oplus V_5.
\]
Entonces
\[
V_{11}=\ell_K\oplus V_{10},
\qquad
\dim V_{10}=10.
\]
El proyector sobre la pantalla de dimensión diez es
\[
P_{10}=P_{11}-\frac{u_Ku_K^{\mathsf T}}{\|u_K\|^2}.
\]
\end{theorem}

\begin{proof}
Por el \cref{lem:norma-uk}, \(u_K\neq0\) y, por definición, \(u_K\in V_3\).
De aquí
\[
V_3=\ell_K\oplus(V_3\cap\ell_K^\perp).
\]
Al añadir los sumandos ortogonales \(V_{3'}\) y \(V_5\), el
\cref{thm:descomposicion-a5} da
\[
V_{11}=\ell_K\oplus
\bigl((V_3\cap\ell_K^\perp)\oplus V_{3'}\oplus V_5\bigr).
\]
La dimensión del segundo sumando es \(2+3+5=10\). La fórmula de \(P_{10}\)
resta a \(P_{11}\) el proyector ortogonal de rango uno sobre \(\ell_K\), de
modo que su imagen es precisamente \(V_{10}\).
\end{proof}

\begin{corollary}[Descomposición exacta \(1+3+7\)]
\label{cor:descomposicion-1-3-7}
Sea
\[
W_7(K):=(V_3\cap\ell_K^\perp)\oplus V_5.
\]
Entonces
\[
\boxed{
V_{11}=\ell_K\oplus V_{3'}\oplus W_7(K),
\qquad
\dim(\ell_K,V_{3'},W_7)=(1,3,7),}
\]
y
\[
\boxed{
V_{10}=V_{3'}\oplus W_7(K),
\qquad 10=3+7.}
\]
\end{corollary}

\begin{proof}
\(\ell_K\subset V_3\) es una recta, de modo que
\(\dim(V_3\cap\ell_K^\perp)=2\). Por tanto,
\(\dim W_7=2+5=7\), y las dos descomposiciones siguen del
\cref{thm:once-diez}.
\end{proof}

\begin{remark}[Condición adicional para la lectura \(1+3+6\)]
\label{rem:interfaz-1-3-6}
Las descomposiciones exactas \(1+3+7\) y \(3+7\) no seleccionan por sí solas
un tiempo, un espacio extendido tridimensional ni un sector compacto de rango
seis. Una realización física debe elegir una recta
\(\ell_t\subset W_7(K)\), demostrar
\[
W_7(K)=\ell_t\oplus W_6,
\]
construir una retícula de rango pleno
\(\Lambda_6\subset W_6\), estable bajo el transporte, interpretar
\(W_6/\Lambda_6\) como el sector compacto y exhibir un entrelazador entre el
lector temporal del TRIT, \(W_6\) y esa retícula. El mismo transporte debe
preservar separadamente las tres direcciones extendidas del sumando
\(V_{3'}\), sin mezclarlas con \(\ell_t\) ni con el cociente compacto. Sólo
bajo esas hipótesis se obtienen las descomposiciones
\[
\boxed{
V_{11}=\ell_K\oplus\ell_t\oplus V_{3'}\oplus W_6,
\qquad 11=1+1+3+6,}
\]
\[
\boxed{
V_{10}=\ell_t\oplus V_{3'}\oplus W_6,
\qquad 10=1+3+6.}
\]
La ausencia de una elección natural de \(\ell_t\), de estabilidad de
\(\Lambda_6\), de rango pleno, de separación de las tres direcciones
extendidas o del entrelazador refutaría esta realización, sin alterar el
corolario algebraico \(1+3+7\).
\end{remark}

La reducción no es \(A_5\)-equivariante en ausencia del dato marcado. El
submódulo \(V_3\) es irreducible y no contiene una recta \(A_5\)-invariante;
la dirección \(\ell_K\) aparece después de introducir \(K\). La secuencia

\[
12\longrightarrow11\longrightarrow10
\]

posee así dos naturalezas distintas: el paso \(12\to11\) es canónico para
la acción de permutación, mientras que \(11\to10\) es una polarización
relativa al vector \(K\) y a la correspondencia posición--clase lateral declarada.

\section{Reducción isotrópica de rango \(26\to24\)}
\label{sec:reduccion-26-24}

Sea \(U\) el plano hiperbólico entero con base \(e_+,e_-\) y matriz de Gram

\[
\begin{pmatrix}0&1\\1&0\end{pmatrix}.
\]

Es un retículo par, unimodular y de firma \((1,1)\). El retículo

\[
L_{26}:=\Lambda_{24}\oplus U
\]

es par, unimodular y de firma \((25,1)\); por la clasificación de retículos
indefinidos pares unimodulares se identifica con \(II_{25,1}\)
\cite{ConwaySloane1999}.

El adjetivo \emph{lorentziano} se emplea en sentido algebraico: indica la
presencia de una dirección negativa en la forma bilineal. La terminología
procede de la geometría de Lorentz y Minkowski
\cite{Lorentz1904,Minkowski1909}; no identifica por sí sola el retículo con
un espacio-tiempo físico.

\begin{theorem}[Reducción nula]
\label{thm:reduccion-nula}
Para el vector isotrópico primitivo \(e_+\in U\), existe un isomorfismo
canónico, una vez fijada la descomposición \(L_{26}=\Lambda_{24}\oplus U\),
\[
e_+^\perp/\mathbb Ze_+\cong\Lambda_{24}.
\]
\end{theorem}

\begin{proof}
Todo elemento de \(L_{26}\) se escribe de manera única como
\(x=\lambda+ae_++be_-\), con
\(\lambda\in\Lambda_{24}\) y \(a,b\in\mathbb Z\). Puesto que la suma es
ortogonal y \(\langle e_-,e_+\rangle=1\),
\[
\langle x,e_+\rangle=b.
\]
Por tanto,
\[
e_+^\perp=\Lambda_{24}\oplus\mathbb Ze_+.
\]
El cociente por \(\mathbb Ze_+\) elimina el segundo sumando y conserva el
primero, lo que define el isomorfismo afirmado.
\end{proof}

La dimensión disminuye en dos porque la condición de ortogonalidad elimina
una dirección y el cociente elimina la dirección nula restante. No se trata
de una proyección ortogonal sobre un complemento positivo: una recta
isotrópica está contenida en su propio ortogonal. Ésa es la razón algebraica
por la que el procedimiento correcto es \(e_+^\perp/\mathbb Ze_+\).

La matriz del plano \(U\) coincide formalmente con la matriz que intercambia
los dos sectores del transporte bidireccional. Adoptarla como forma de Gram
establece una realización reticular concreta de ese involutor. La reducción
\cref{thm:reduccion-nula} es entonces exacta dentro de dicha realización;
la coincidencia matricial, por sí sola, no identifica el retículo con un
espacio-tiempo físico.

Este ambiente es también el empleado en la construcción de Borcherds: el
sector de carga central veinticuatro se combina con el plano lorentziano de
rango dos antes de la reducción que conduce al álgebra de Lie monstruosa
\cite{Borcherds1992}. La relación pertinente aquí es el cociente reticular
explícito; la teoría de álgebras de operadores de vértice sigue el corredor
clásico descrito anteriormente.


\section{Correspondencia icosaédrica de McKay y tipo \texorpdfstring{\(E_8\)}{E8}}
\label{sec:tres-ochos}

Tres apariciones del entero ocho intervienen cerca de esta frontera:

\[
8_{\rm oct}=|O_H|,
\qquad
V_8:=V_{3'}\oplus V_5,
\qquad
\operatorname{rank}(E_8)=8.
\]

La primera es el cardinal de una octada del sistema de Witt; la segunda es
un módulo real de \(A_5\); la tercera es el rango de un sistema de raíces. No
son instancias de un único objeto. El puente clásico relevante pasa por el
grupo icosaédrico binario:

\[
A_5\longleftarrow2A_5\subset SU(2)
\xrightarrow{\text{McKay}}\widetilde E_8
\]

\cite{McKay1980}. En la correspondencia de McKay, los vértices del diagrama
afín codifican representaciones irreducibles de \(2A_5\), y las aristas
describen la descomposición del producto tensorial con la representación
bidimensional natural. Esta construcción explica por qué el icosaedro y
\(E_8\) están vinculados; no identifica una octada de Witt con el módulo
\(V_8\), ni uno de éstos con el retículo \(E_8\).

\begin{proposition}[Separación categorial de tres invariantes de valor ocho]
Ninguna igualdad entre los tres invariantes anteriores define por sí misma
un entrelazador. Para relacionar dos de ellos se necesita, respectivamente,
un mapa de conjuntos de incidencia a módulos, de módulos de \(A_5\) a
módulos de \(2A_5\), o de representaciones a datos reticulares.
\end{proposition}

\begin{proof}
El cardinal de un conjunto, la dimensión de un espacio vectorial y el rango
de un retículo son invariantes en categorías diferentes. Un isomorfismo en
una de ellas no está definido para los objetos de las otras. La
correspondencia de McKay proporciona uno de los funtores intermedios, pero no
los restantes; por ello la coincidencia del entero ocho es insuficiente.
\end{proof}

\section{Reducción dimensional y contexto de la teoría M}
\label{sec:interfaz-teoria-m}

Las construcciones anteriores producen dos apariciones exactas de once. La
primera es la longitud del código perfecto \(G_{11}\); la segunda es la
dimensión del módulo activo \(V_{11}=\mathbf1^\perp\). Estos objetos no se
identifican: uno está definido sobre \(\mathbb F_3\) y el otro sobre
\(\mathbb R\). La reducción a diez pertenece exclusivamente al segundo y
depende de la recta marcada \(\ell_K\).

La supergravedad en once dimensiones y la red de dualidades que condujo a la
formulación de la teoría M aportan el contexto físico en el que una dirección
undecadimensional puede compactificarse y producir teorías de diez
dimensiones \cite{CremmerJuliaScherk1978,HullTownsend1995,Witten1995,Townsend1995}.
La correspondencia dimensional sugerida por la construcción HMT se organiza
en el siguiente diccionario:

\[
\begin{array}{>{\raggedright\arraybackslash}p{0.42\textwidth}
                >{\raggedright\arraybackslash}p{0.46\textwidth}}
\toprule
\textbf{Objeto algebraico} & \textbf{Dato requerido en una realización física}\\
\midrule
\(V_{11}=\ell_K\oplus V_{10}\) &
una dirección compacta y una pantalla de diez dimensiones con un mapa de
campos definido\\
\((m,w,R_0)\mapsto(w,m,R_0^{-1})\) &
sectores de momento y enrollamiento, radio físico y escala de cuerda\\
\(\Lambda_{24}\oplus U\cong II_{25,1}\) &
una interpretación dinámica del sector reticular lorentziano\\
\(C_W\to G_{11}\) y \(C_W\to\Lambda_{24}\) &
un entrelazador que preserve incidencia, acción y espectro\\
\bottomrule
\end{array}
\]

En la supergravedad undecadimensional, una realización completa debe incluir,
al menos, el campo métrico \(g_{MN}\), la tres-forma \(C_{MNP}\), el
gravitino \(\psi_M\), su acción y sus transformaciones de supersimetría. Para
que la pantalla \(V_{11}=\ell_K\oplus V_{10}\) represente una compactificación
física, hace falta una aplicación que envíe esos campos a datos sobre
\(\ell_K\) y \(V_{10}\), y que preserve la dinámica y el límite de baja
energía. Ninguno de esos campos entra como hipótesis en los
\cref{thm:g11-perfecto,thm:once-diez,thm:reduccion-nula,thm:dualidad-visible-memoria}.

Por consiguiente, la interfaz posee dos niveles bien diferenciados. En el
nivel algebraico quedan demostradas la perfección del código, la
descomposición \(12=1+3+3'+5\), la polarización \(11=1+10\), la reducción
isotrópica \(26\to24\) y la involución de radio. En el nivel físico, estas
estructuras proporcionan un dominio y unas ecuaciones de compatibilidad para
construir un diccionario con teoría M. La segunda tarea añade campos y
dinámica; no modifica ni se usa para demostrar el primer nivel.

\Needspace{24\baselineskip}
\section{Diagrama de cocientes y reducciones de rango}
\label{sec:sintesis-pantallas}

El diagrama de rangos puede resumirse sin perder sus tipos:

\[
\begin{array}{ccccc}
&&C_W=[12,6,6]_3&&\\[1mm]
&\swarrow&&\searrow&\\[-1mm]
G_{11}\subset\mathbb F_3^{11}&&&&N(A_2^{12})\longrightarrow\Lambda_{24}\\[2mm]
&&V_{12}=\mathbb R^{12}&&\\[-1mm]
&&\downarrow P_{11}&&\\[-1mm]
&&V_{11}=V_3\oplus V_{3'}\oplus V_5&&\\[-1mm]
&&\downarrow P_{10}(K)&&\\[-1mm]
&&V_{10}&&\\[2mm]
&&II_{25,1}=\Lambda_{24}\oplus U&&\\[-1mm]
&&\downarrow\ e_+^\perp/\mathbb Ze_+&&\\[-1mm]
&&\Lambda_{24}.&&
\end{array}
\]

La flecha \(P_{11}\) elimina el modo uniforme; \(P_{10}(K)\) elimina la
dirección seleccionada por \(K\); la reducción nula elimina una pareja
isotrópica mediante ortogonal y cociente. La involución
\(\mathsf T\) actúa en un dominio reticular ampliado y relaciona radios
recíprocos. Éstos son los mecanismos matemáticos que sostienen la interfaz
dimensional. Su formulación explícita permite distinguir con precisión los
teoremas de rango de las interpretaciones físicas que pueden edificarse
sobre ellos.
